\pdfbookmark[0]{Streszczenie}{streszczenie.1}
%\phantomsection
%\addcontentsline{toc}{chapter}{Streszczenie}
%%% Poni¿sze zosta³o niewykorzystane (tj. zrezygnowano z utworzenia nienumerowanego rozdzia³u na abstrakt)
%%%\begingroup
%%%\setlength\beforechapskip{48pt} % z jakiegoœ powodu by³a maleñka ró¿nica w po³o¿eniu nag³ówka rozdzia³u numerowanego i nienumerowanego
%%%\chapter*{\centering Abstrakt}
%%%\endgroup
%%%\label{sec:abstrakt}
%%%Lorem ipsum dolor sit amet eleifend et, congue arcu. Morbi tellus sit amet, massa. Vivamus est id risus. Sed sit amet, libero. Aenean ac ipsum. Mauris vel lectus. 
%%%
%%%Nam id nulla a adipiscing tortor, dictum ut, lobortis urna. Donec non dui. Cras tempus orci ipsum, molestie quis, lacinia varius nunc, rhoncus purus, consectetuer congue risus. 
%\mbox{}\vspace{2cm} % mo¿na przesun¹æ, w zale¿noœci od d³ugoœci streszczenia
\begin{abstract}
Praca dotyczy zaprojektowania i walidacji architektury \emph{Shadow Deploy} dla krytycznych funkcji 5G Core uruchamianych jako \emph{Cloud Native Network Functions} w środowisku Kubernetes. Głównym celem było przygotowanie bezpiecznego mechanizmu wdrażania nowych wersji funkcji SMF bez ryzyka wpływu na ruch produkcyjny abonentów.

W pracy zdefiniowano dwa filary rozwiązania:
\begin{itemize}
    \item kontrolę ruchu i duplikację żądań z użyciem Istio/Envoy (\emph{traffic mirroring}),
    \item izolację stanu danych przez rozdzielenie bazy produkcyjnej UDR od bazy testowej \texttt{udr-shadow}.
\end{itemize}

\vspace{2em}

Warstwa Service Mesh została opisana w oparciu o obiekty \texttt{DestinationRule} i \texttt{VirtualService}, które kierują 100\% ruchu do wersji stabilnej SMF v1 oraz kopiują część żądań do wersji shadow (SMF v2). Wymusza to rzeczywistą walidację nowego kodu przy zachowaniu pełnej transparentności dla użytkownika końcowego, ponieważ odpowiedzi shadow nie wracają do klienta.

Kluczowym elementem pracy jest model oceny jakości oparty na metrykach z Envoy i Prometheus. Zastosowano progi SLO dla opóźnień i błędów HTTP oraz scenariusz oceny eksperymentu z wykorzystaniem Iter8. Uzupełniająco przedstawiono konfigurację narzędzia Diffy do porównania odpowiedzi wersji stabilnej i shadow.

Wyniki pracy potwierdzają, że architektura Shadow Deploy może znacząco obniżyć ryzyko wdrożeń w 5G Core: pozwala wykrywać regresje wydajnościowe przed przełączeniem ruchu, zwiększa bezpieczeństwo zmian oraz wspiera procesy DevOps/SRE dla środowisk telekomunikacyjnych o wysokich wymaganiach niezawodności.
\end{abstract}
\vspace{2em}
\mykeywords
\pagebreak
{
\selectlanguage{english}
\begin{abstract}
This thesis presents a Shadow Deployment architecture for critical 5G Core functions executed as Cloud Native Network Functions on Kubernetes. The objective is to validate new SMF versions against real production traffic patterns without affecting subscribers.

The proposed architecture combines Istio/Envoy traffic mirroring with strict data-state isolation. In the runtime model, 100\% of user traffic is served by the stable SMF v1, while mirrored requests are sent to SMF v2 for non-intrusive verification. The v2 instance is hard-wired to an isolated \texttt{udr-shadow} datastore to prevent accidental writes to production subscriber data.

The evaluation framework is based on service-mesh telemetry, Prometheus metrics, and Iter8 quality gates. Service-level objectives are defined for latency and HTTP error ratio, enabling automated pass/fail assessment of the shadow version. Additionally, an optional Diffy-based response comparator is included for payload-level behavioral analysis.

The work shows that Shadow Deployment can reduce rollout risk in 5G Core environments, improve release confidence, and provide a measurable validation path aligned with DevOps/SRE practices and carrier-grade reliability requirements.
\end{abstract}
\vspace{2em}
\mykeywords
}
